% math0 manuscript. Build with: latexmk -pdf paper.tex
\documentclass[11pt]{amsart}
\usepackage[margin=1.1in]{geometry}
\usepackage{amsmath,amssymb,amsthm,mathtools}
\usepackage{url}
\usepackage{hyperref}
\usepackage{xcolor}
\usepackage{listings}
\usepackage{newunicodechar}

\lstdefinelanguage{lean}{
  keywords={theorem,lemma,def,abbrev,example,fun,by,have,show,from,let,in,if,then,else,match,with,
    namespace,end,open,import,section,variable,instance,structure,where,calc,obtain,intro,exact,
    refine,rcases,rintro,simp,omega,decide,norm_num,linarith,nlinarith,positivity,induction,cases,
    noncomputable,private,rw,apply,interval_cases},
  sensitive=true, comment=[l]{--}, morecomment=[s]{/-}{-/}, string=[b]"}
\lstset{language=lean, basicstyle=\ttfamily\footnotesize, columns=fullflexible, keepspaces=true, breaklines=true,
  keywordstyle=\color{blue!60!black}\bfseries, commentstyle=\color{gray}, frame=single,
  extendedchars=true, inputencoding=utf8,
  literate={ℕ}{{$\mathbb{N}$}}1 {∀}{{$\forall$}}1 {∃}{{$\exists$}}1 {∈}{{$\in$}}1 {→}{{$\to$}}1
    {↔}{{$\leftrightarrow$}}1 {∧}{{$\wedge$}}1 {∨}{{$\vee$}}1 {≤}{{$\le$}}1 {≥}{{$\ge$}}1
    {¬}{{$\neg$}}1 {·}{{$\cdot$}}1 {←}{{$\leftarrow$}}1}
% Unicode commonly found in Lean sources.
\newunicodechar{ℕ}{\ensuremath{\mathbb{N}}}
\newunicodechar{ℤ}{\ensuremath{\mathbb{Z}}}
\newunicodechar{ℚ}{\ensuremath{\mathbb{Q}}}
\newunicodechar{ℝ}{\ensuremath{\mathbb{R}}}
\newunicodechar{ℂ}{\ensuremath{\mathbb{C}}}
\newunicodechar{∀}{\ensuremath{\forall}}
\newunicodechar{∃}{\ensuremath{\exists}}
\newunicodechar{→}{\ensuremath{\to}}
\newunicodechar{←}{\ensuremath{\leftarrow}}
\newunicodechar{↔}{\ensuremath{\leftrightarrow}}
\newunicodechar{↦}{\ensuremath{\mapsto}}
\newunicodechar{¬}{\ensuremath{\neg}}
\newunicodechar{∧}{\ensuremath{\wedge}}
\newunicodechar{∨}{\ensuremath{\vee}}
\newunicodechar{≤}{\ensuremath{\le}}
\newunicodechar{≥}{\ensuremath{\ge}}
\newunicodechar{≠}{\ensuremath{\ne}}
\newunicodechar{∈}{\ensuremath{\in}}
\newunicodechar{∉}{\ensuremath{\notin}}
\newunicodechar{⊆}{\ensuremath{\subseteq}}
\newunicodechar{∩}{\ensuremath{\cap}}
\newunicodechar{∪}{\ensuremath{\cup}}
\newunicodechar{∑}{\ensuremath{\sum}}
\newunicodechar{∏}{\ensuremath{\prod}}
\newunicodechar{⟨}{\ensuremath{\langle}}
\newunicodechar{⟩}{\ensuremath{\rangle}}
\newunicodechar{·}{\ensuremath{\cdot}}
\newunicodechar{∘}{\ensuremath{\circ}}
\newunicodechar{⁻}{\ensuremath{^{-}}}
\newunicodechar{¹}{\ensuremath{^{1}}}
\newunicodechar{²}{\ensuremath{^{2}}}
\newunicodechar{₀}{\ensuremath{_{0}}}
\newunicodechar{₁}{\ensuremath{_{1}}}
\newunicodechar{₂}{\ensuremath{_{2}}}
\newunicodechar{α}{\ensuremath{\alpha}}
\newunicodechar{β}{\ensuremath{\beta}}
\newunicodechar{ε}{\ensuremath{\varepsilon}}
\newunicodechar{δ}{\ensuremath{\delta}}
\newunicodechar{π}{\ensuremath{\pi}}
\newunicodechar{σ}{\ensuremath{\sigma}}
\newunicodechar{φ}{\ensuremath{\varphi}}
\newunicodechar{∞}{\ensuremath{\infty}}
\newunicodechar{⌊}{\ensuremath{\lfloor}}
\newunicodechar{⌋}{\ensuremath{\rfloor}}
\newunicodechar{⌈}{\ensuremath{\lceil}}
\newunicodechar{⌉}{\ensuremath{\rceil}}
\newunicodechar{∣}{\ensuremath{\mid}}
\newunicodechar{≡}{\ensuremath{\equiv}}
\newunicodechar{×}{\ensuremath{\times}}
\newunicodechar{⧸}{\ensuremath{/}}
\newunicodechar{‖}{\ensuremath{\|}}
\newunicodechar{√}{\ensuremath{\surd}}
\newunicodechar{∅}{\ensuremath{\emptyset}}
\newunicodechar{⦃}{\ensuremath{\{\!|}}
\newunicodechar{⦄}{\ensuremath{|\!\}}}
\newunicodechar{𝓝}{\ensuremath{\mathcal{N}}}
\newunicodechar{≫}{\ensuremath{\gg}}
\newunicodechar{≪}{\ensuremath{\ll}}

\theoremstyle{plain}
\newtheorem{theorem}{Theorem}
\newtheorem{lemma}[theorem]{Lemma}
\newtheorem{proposition}[theorem]{Proposition}
\newtheorem{corollary}[theorem]{Corollary}
\theoremstyle{definition}
\newtheorem{definition}[theorem]{Definition}
\newtheorem{remark}[theorem]{Remark}

\emergencystretch=3em
% Lean identifiers in running text (breakable at dots and underscores).
\expandafter\def\expandafter\UrlBreaks\expandafter{\UrlBreaks\do\_\do\.}
\DeclareUrlCommand\decl{\urlstyle{tt}}
\newcommand{\N}{\mathbb{N}}
\newcommand{\rep}{\operatorname{rep}}
\newcommand{\dprod}{\operatorname{p}}

\title[Repdigits with repdigit digital product (OEIS A201060)]{A repdigit with at least four digits whose digital product is a repdigit is a repunit: the conjecture in OEIS A201060}
\author{math0}
\address{Produced by Claude (an AI model developed by Anthropic) in the math0 project}
\date{October 20, 2026}

\begin{document}
\begin{abstract}
Let $a(1)<a(2)<\cdots$ be the natural numbers $m$ such that $m$ and the product of the decimal digits of $m$ are both repdigits; this is OEIS A201060, which begins $0,1,\dots,9,11,22,33,111,222,1111,11111,\dots$.
We prove the conjecture recorded in the entry by J.~Krizek (2011) that $a(n)$ is a repunit for every $n\ge16$; equivalently, $d^k$ is not a repdigit when $2\le d\le9$ and $k\ge4$.
The result is not new in substance: it follows at once from Obl\'ath's theorem (1956) on perfect powers with all digits equal to some $a\ne1$, together with the result of Bugeaud and Mignotte (1999) that no decimal repunit with at least two digits is a perfect power; the OEIS entry nevertheless still lists it as a conjecture.
We give a short elementary proof, using divisibility by $16$ for even $d$ and residues modulo $100$ for odd $d$, and, to our knowledge, the first formal proof: it is verified in Lean~4 with Mathlib against a statement that math0 formalized from the OEIS entry, reviewed for fidelity and frozen before any proof attempt, whose definitions are checked by the kernel against all $31$ listed terms.
\end{abstract}
\maketitle

\section{Introduction}

A \emph{repdigit} is a natural number whose decimal digits are all equal, such as $7$, $44$ or $333$; following A010785 \cite{A010785}, $0$ is a repdigit. The \emph{repunits} are the repdigits $R_k=(10^k-1)/9$ whose digits are all $1$ (A002275 \cite{A002275}), and the \emph{digital product} $\dprod(m)$ of $m$ is the product of its decimal digits (A007954 \cite{A007954}).
The OEIS entry A201060 \cite{A201060}, created by Jaroslav Krizek on November~26, 2011, lists the repdigits with repdigit digital product:
\[
0,1,2,3,4,5,6,7,8,9,11,22,33,111,222,1111,11111,111111,\dots,1111111111111111111
\]
(offset $1$, $31$ listed terms). Its comment section contains the conjecture that $a(n)$ is a repunit for $n\ge16$ (quoted in Section~\ref{sec:statement}). We prove it (Theorem~\ref{thm:main}).

A repdigit with $k\ge1$ digits all equal to $d\in\{1,\dots,9\}$ has digital product $d^k$. Every repunit is therefore a term, and since $a(16)=1111$, the conjecture says that no repdigit with $k\ge4$ digits all equal to some $d\in\{2,\dots,9\}$ is a term, that is,
\begin{equation}\label{eq:reform}
  d^k \text{ is not a repdigit whenever } 2\le d\le 9 \text{ and } k\ge 4 .
\end{equation}

\subsection*{Prior work and status}
Statement~\eqref{eq:reform} is a special case of classical results on perfect powers with identical digits, and we do not claim it as new.
Obl\'ath \cite{Obl56} proved that the only perfect powers all of whose decimal digits are equal to some fixed digit $a\ne1$ are $4$, $8$ and $9$; we rely on the restatement of this result by Kihel and Luca \cite[Introduction]{KL05} (see also Gica and Panaitopol \cite{GP03}, who describe Obl\'ath's work as almost entirely solving the problem of finding all perfect powers with equal digits).
The remaining case, of repunits, is covered by the work of Bugeaud and Mignotte \cite{BM99}: as Kihel and Luca note \cite[Introduction]{KL05}, thanks to their results the only solutions of $(x^n-1)/(x-1)=y^q$ in integers $n\ge3$, $y\ge2$, $q\ge2$ with base $2\le x\le10$ are $(3^5-1)/2=11^2$ and $(7^4-1)/6=20^2$ (the list in \cite{KL05} also contains $(18^3-1)/17=7^3$, whose base lies outside this range). None of them has $x=10$, so no decimal repunit $R_n$ with $n\ge3$ is a perfect power; $R_2=11$ is not one either.
Together, these results say that no repdigit $N\ge10$ is a perfect power. For $2\le d\le9$ and $k\ge4$ the number $d^k\ge16$ is a perfect power, and \eqref{eq:reform} follows.
The OEIS entry, in revision \#11 of March~30, 2012, still records the statement as a conjecture, and neither it nor the two entries it cross-references (A201061 and A201064, also by Krizek) mention a proof. The math0 referees found no source that discusses A201060 by name, and the statement is not among those of the formal-conjectures repository \cite{FC}.
We therefore classify our result as the first formal proof, to our knowledge, of a known result.

\subsection*{Method}
Because the base $d$ of the power in \eqref{eq:reform} is a single digit, the deep results above are not needed, and our proof uses only congruences. A positive repdigit is $eR_j$ with $1\le e\le9$ and $R_j$ odd, so it is never divisible by $16$; this excludes even $d$. A repdigit $N\ge10$ ends in the two digits $ee$, so $N\bmod 100=11e$ is a multiple of $11$, whereas the powers of $3$, $5$, $7$ and $9$ modulo $100$ never are; this excludes odd $d$.
The Lean proof follows exactly this route (Section~\ref{sec:proof}).

\section{Statement}\label{sec:statement}

\begin{definition}\label{def:main}
For $m\in\N$ let $\delta(m)$ be the (possibly empty) sequence of decimal digits of $m$; thus $\delta(0)$ is empty. Call $m$ a \emph{repdigit} if all entries of $\delta(m)$ are equal, and a \emph{repunit} if all entries of $\delta(m)$ equal $1$. Let $\dprod(0)=0$ and, for $m\ge1$, let $\dprod(m)$ be the product of the entries of $\delta(m)$. Let $S=\{m\in\N: m \text{ and } \dprod(m) \text{ are repdigits}\}$, and let $a(1)<a(2)<\cdots$ be the elements of $S$ in increasing order.
\end{definition}

The set $S$ is infinite, since it contains every repunit $R_k$ ($k\ge1$), whose digital product is $1$; so $a(n)$ is defined for every $n\ge1$.
The comment section of A201060 \cite{A201060} contains, verbatim:
\begin{quote}
{\small\texttt{Conjecture: a(n) is repunit for n >= 16.}}
\end{quote}
The entry gives no date for this comment; the entry itself is credited to Jaroslav Krizek, November~26, 2011.

\begin{theorem}\label{thm:main}
For every $n\ge16$, $a(n)$ is a repunit.
\end{theorem}

\begin{remark}\label{rem:closed}
Combined with the listed data, Theorem~\ref{thm:main} determines the sequence: the terms below $1111$ are $a(1),\dots,a(15)=0,1,\dots,9,11,22,33,111,222$, and every repunit $R_k$ with $k\ge4$ is a term, so $a(n)=R_{n-12}=(10^{n-12}-1)/9$ for all $n\ge16$. This closed form is not stated as a separate Lean declaration; its ingredients are Theorem~\ref{thm:main}, Lemma~\ref{lem:seq} and the fact that every repunit is a term.
\end{remark}

\subsection*{Lean formalization}
The statement was formalized by math0 in the file \path{lean/Math0/Statements/OEIS/A201060.lean}, and was reviewed for fidelity and frozen (SHA-256 \texttt{902fc983\ldots 2af72b}) before any proof attempt. Its definitions and the statement read as follows (verbatim; the copyright header, the import, the module docstring, section headers, docstrings, blank lines, the two decidability instances, the auxiliary definitions \decl{rep}, \decl{repList} and \decl{listed}, and all theorems are omitted):
\begin{lstlisting}
namespace Math0.Statements.A201060
def IsRepdigit (m : ℕ) : Prop := ∀ x ∈ Nat.digits 10 m, ∀ y ∈ Nat.digits 10 m, x = y
def digitProd (m : ℕ) : ℕ := if m = 0 then 0 else (Nat.digits 10 m).prod
def IsTerm (m : ℕ) : Prop := IsRepdigit m ∧ IsRepdigit (digitProd m)
noncomputable def a (n : ℕ) : ℕ := Nat.nth IsTerm (n - 1)
def IsRepunit (m : ℕ) : Prop := ∀ d ∈ Nat.digits 10 m, d = 1
def statement : Prop := ∀ n : ℕ, 16 ≤ n → IsRepunit (a n)
end Math0.Statements.A201060
\end{lstlisting}
Here \texttt{Nat.digits 10 m} is the list $\delta(m)$ of decimal digits of $m$, least significant first, with \texttt{Nat.digits 10 0 = []}; so \decl{IsRepdigit}, \decl{IsRepunit} and \decl{digitProd} are the notions of Definition~\ref{def:main}, and $0$ is a repdigit, as in A010785 and in the entry ($a(1)=0$). The convention $\dprod(0)=0$ follows A007954; the empty product $1$ would give the same sequence, since $0$ and $1$ are both repdigits. \decl{IsRepunit} also holds for $0$ (as A002275 includes $R_0=0$), which is irrelevant here because $a(n)\ge a(16)=1111$ for $n\ge16$.
\texttt{Nat.nth p k} is the $k$-th element, counting from $0$, of $\{m: p(m)\}$ in increasing order, so \decl{a} is the sequence of Definition~\ref{def:main} with the entry's offset $1$. When the set is infinite, \texttt{Nat.nth p} is strictly increasing and enumerates the whole set; the statement file proves that $S$ is infinite. The statement keeps the entry's range $n\ge16$ and adds no hypotheses.

To compare the definitions with the OEIS data, the statement file introduces $\rep(d,k)=\sum_{i<k}d\cdot10^i$ (in Lean \texttt{rep d k = Nat.ofDigits 10 (List.replicate k d)}), the number with $k$ digits all equal to $d$, and proves, among others, the following lemmas (statements shown verbatim, docstrings and proofs omitted):
\begin{lstlisting}
theorem digits_rep {d : ℕ} (hd1 : 1 ≤ d) (hd9 : d ≤ 9) (k : ℕ) :
    Nat.digits 10 (rep d k) = List.replicate k d
theorem isRepdigit_iff_mem_repList {L m : ℕ} (hm : m < 10 ^ L) :
    IsRepdigit m ↔ m ∈ repList L
theorem infinite_isTerm : (Set.ofPred IsTerm).Infinite
theorem a_data : (List.range 31).map (fun i => a (i + 1)) =
    [0, 1, 2, 3, 4, 5, 6, 7, 8, 9, 11, 22, 33, 111, 222, 1111, 11111, 111111, 1111111, 11111111,
      111111111, 1111111111, 11111111111, 111111111111, 1111111111111, 11111111111111,
      111111111111111, 1111111111111111, 11111111111111111, 111111111111111111,
      1111111111111111111]
\end{lstlisting}
Here \texttt{repList L} is the explicit list consisting of $0$ and the numbers $\rep(d,k)$ with $1\le d\le9$, $1\le k\le L$; so \decl{isRepdigit_iff_mem_repList} says that the repdigits below $10^L$ are exactly $0$ and these numbers.
The theorem \decl{a_data} checks all $31$ listed terms. Its proof does not enumerate the numbers below $10^{19}$. Instead, the kernel filters the $1+9\cdot19=172$ elements of \texttt{repList 19} by the condition ``$\dprod(m)$ is a repdigit'' and checks that the result is the list of the $31$ listed terms, that every element of \texttt{repList 19} is below $10^{19}$, and that the listed terms are strictly increasing (three proofs by \texttt{decide +kernel}). By \decl{isRepdigit_iff_mem_repList} with $L=19$, the terms below $10^{19}$ and the listed terms then have the same elements; both lists are strictly increasing, so they are equal, and a general lemma (\decl{map_nth_eq_of_filter}) converts the statement about the sorted list of terms below $10^{19}$ into the values \texttt{Nat.nth IsTerm i} for $i<31$.

\section{Proof}\label{sec:proof}

The proof follows the Lean file \path{lean/Math0/Frontier/OEIS/A201060.lean}; the name of the corresponding Lean declaration is given with each step. We write $\rep(e,j)$ as above and $R_j=\rep(1,j)$.

\begin{lemma}[\decl{digits_rep}, \decl{rep_succ}, \decl{rep_eq_mul}, \decl{odd_rep_one}]\label{lem:rep}
For all $e,j\in\N$:
\begin{enumerate}
\item if $1\le e\le9$, the decimal digits of $\rep(e,j)$ are $j$ copies of $e$;
\item $\rep(e,j+1)=e+10\,\rep(e,j)$;
\item $\rep(e,j)=e\,R_j$;
\item $R_{j+1}$ is odd.
\end{enumerate}
\end{lemma}

\begin{proof}
(1) holds because the digit list $(e,\dots,e)$ consists of digits less than $10$ and has a nonzero last (leading) entry, so it is the digit list of the number it represents. (2) is the definition of $\rep$ via \texttt{Nat.ofDigits}. (3) follows from (2) by induction on $j$, with $\rep(e,0)=0$. By (2), $R_{j+1}=1+10R_j=2\cdot5R_j+1$, which gives (4).
\end{proof}

\begin{lemma}[\decl{isRepdigit_iff_mem_repList}, \decl{repdigit_eq}]\label{lem:shape}
Every repdigit $N\ge1$ equals $\rep(e,j)$ for some $1\le e\le9$ and $j\ge1$.
\end{lemma}

\begin{proof}
Since $N\ge1$, its digit list is nonempty; let $e$ be its first entry and $j\ge1$ its length. All entries equal $e$, so $N=\rep(e,j)$. Moreover $e\le9$ because $e$ is a digit, and $e\ge1$ because the last (leading) digit of a positive number is nonzero. In Lean this is the forward direction of \decl{isRepdigit_iff_mem_repList}, applied with $L=N$ (using $N<10^N$).
\end{proof}

\begin{lemma}[\decl{repdigit_not_dvd_sixteen}]\label{lem:16}
No repdigit $N\ge1$ is divisible by $16$.
\end{lemma}

\begin{proof}
By Lemmas~\ref{lem:shape} and~\ref{lem:rep}(3), $N=e\,R_j$ with $1\le e\le9$ and $j\ge1$. By Lemma~\ref{lem:rep}(4), $R_j$ is odd, so it is coprime to $2$ and hence to $16=2^4$. If $16\mid eR_j$, then $16\mid e$, which forces $16\le e$ since $e\ge1$; this contradicts $e\le9$.
\end{proof}

\begin{lemma}[\decl{repdigit_mod_hundred}]\label{lem:100}
If $N\ge10$ is a repdigit, then $N\bmod 100$ is divisible by $11$.
\end{lemma}

\begin{proof}
By Lemma~\ref{lem:shape}, $N=\rep(e,j)$ with $1\le e\le9$ and $j\ge1$. If $j=1$ then $N=e\le9$, which is excluded; so $j=i+2$ for some $i\ge0$, and Lemma~\ref{lem:rep}(2), applied twice, gives
\[
  N=e+10\bigl(e+10\,\rep(e,i)\bigr)=11e+100\,\rep(e,i).
\]
Since $11e\le99$, we get $N\bmod100=11e$, which is divisible by $11$.
\end{proof}

\begin{lemma}[\decl{pow_mod_mem}]\label{lem:orbit}
Let $D\in\N$ and let $T$ be a finite set of natural numbers with $1\in T$ and $(tD)\bmod100\in T$ for every $t\in T$. Then $D^K\bmod 100\in T$ for every $K\ge0$.
\end{lemma}

\begin{proof}
Induction on $K$: $D^0\bmod100=1\in T$, and $D^{K+1}\bmod100=\bigl((D^K\bmod100)\cdot D\bigr)\bmod100$, which lies in $T$ by the hypothesis applied to $t=D^K\bmod100\in T$.
\end{proof}

The two residue sets used are
\begin{align*}
T_3&=\{1,3,9,27,81,43,29,87,61,83,49,47,41,23,69,7,21,63,89,67\},\\
T_5&=\{1,5,25\}.
\end{align*}
$T_3$ is the set of the $20$ residues of $3^0,\dots,3^{19}$ modulo $100$, listed in this order; it is the subgroup of $(\mathbb{Z}/100\mathbb{Z})^\times$ generated by $3$, which contains $7\equiv3^{15}$ and $9=3^2$. In Lean, $T_3$ and $T_5$ are the lists \decl{T3} and \decl{T5}, and the facts needed about them are finite checks:

\begin{lemma}[finite checks, \texttt{decide}]\label{lem:tables}
The following hold.
\begin{enumerate}
\item $1\in T_3$, and $(3t)\bmod100$, $(7t)\bmod100$ and $(9t)\bmod100$ lie in $T_3$ for every $t\in T_3$.
\item $1\in T_5$, and $(5t)\bmod100\in T_5$ for every $t\in T_5$.
\item No element of $T_3\cup T_5$ is divisible by $11$.
\end{enumerate}
\end{lemma}

\begin{proof}
Direct computation. For (3), the residues modulo $11$ of the elements of $T_3$ lie in $\{1,3,4,\dots,10\}$, and $1,5,25\equiv1,5,3\pmod{11}$. In Lean each of the twelve membership and divisibility claims (three for each of $D=3,5,7,9$) is closed by \texttt{decide}, which evaluates the corresponding decision procedure on these explicit lists in the kernel.
\end{proof}

\begin{lemma}[\decl{not_repdigit_even}, \decl{not_repdigit_odd}, \decl{not_repdigit_pow}]\label{lem:pow}
If $2\le D\le9$ and $K\ge4$, then $D^K$ is not a repdigit.
\end{lemma}

\begin{proof}
First, $D^K\ge2^K\ge2^4=16\ge10$. The proof splits into the eight values of $D$ (\texttt{interval\_cases}).

\emph{$D\in\{2,4,6,8\}$.} Since $2\mid D$ and $K\ge4$, we have $16=2^4\mid 2^K\mid D^K$. As $D^K\ge1$, Lemma~\ref{lem:16} shows that $D^K$ is not a repdigit.

\emph{$D\in\{3,5,7,9\}$.} Let $T=T_5$ if $D=5$ and $T=T_3$ otherwise. By Lemmas~\ref{lem:tables}(1),(2) and~\ref{lem:orbit}, $D^K\bmod100\in T$, so by Lemma~\ref{lem:tables}(3), $D^K\bmod100$ is not divisible by $11$. Since $D^K\ge10$, Lemma~\ref{lem:100} shows that $D^K$ is not a repdigit.
\end{proof}

\begin{remark}
The congruence modulo $100$ alone does not suffice for even $D$: for instance $2^{18}=262144\equiv44$ and $2^{19}=524288\equiv88\pmod{100}$, so the last two digits of a power of $2$ can agree. This is why the even case uses divisibility by $16$ instead.
\end{remark}

\begin{lemma}[\decl{isRepunit_of_isTerm}]\label{lem:term}
Every $m\in S$ with $m\ge1111$ is a repunit.
\end{lemma}

\begin{proof}
Since $m$ is a repdigit and $m\ge1$, Lemma~\ref{lem:shape} gives $m=\rep(e,j)$ with $1\le e\le9$ and $j\ge1$. By Lemma~\ref{lem:rep}(1), $m$ has exactly $j$ digits, so $m<10^j$; as $m\ge1111\ge10^3$, this forces $j\ge4$. If $e=1$, all digits of $m$ equal $1$ and $m$ is a repunit. Otherwise $2\le e\le9$, and by Lemma~\ref{lem:rep}(1) the digital product of $m$ is $e^j$, which is a repdigit because $m\in S$. This contradicts Lemma~\ref{lem:pow} with $D=e$, $K=j\ge4$.
\end{proof}

\begin{lemma}[\decl{infinite_isTerm}, \decl{a_data}, \decl{a_sixteen}]\label{lem:seq}
The set $S$ is infinite, and $a(16)=1111$.
\end{lemma}

\begin{proof}
For every $B$, the repunit $R_{B+1}$ lies in $S$ (its digital product is $1$) and satisfies $R_{B+1}\ge10^B>B$; so $S$ is unbounded, hence infinite. The value $a(16)=1111$ is the sixteenth entry of \decl{a_data} (Section~\ref{sec:statement}).
\end{proof}

\begin{proof}[Proof of Theorem~\ref{thm:main}]
Let $n\ge16$. Since $S$ is infinite (Lemma~\ref{lem:seq}), $a(n)$, the element of $S$ with index $n-1$ counting from $0$, lies in $S$ (\texttt{Nat.nth\_mem\_of\_infinite}); and \texttt{Nat.nth} is monotone for an infinite set (\texttt{Nat.nth\_monotone}), so $a(n)\ge a(16)=1111$ because $n-1\ge15$. By Lemma~\ref{lem:term}, $a(n)$ is a repunit.
\end{proof}

\section{Machine verification}\label{sec:lean}

\subsection*{Files and declarations}
The proof is the file \path{lean/Math0/Frontier/OEIS/A201060.lean} of the math0 repository, \url{https://github.com/Foundation-0/math0}. It imports the frozen statement file \path{lean/Math0/Statements/OEIS/A201060.lean} (Section~\ref{sec:statement}) and the audit commands of \path{lean/Math0/Audit.lean}. The main declaration is \decl{Math0.Frontier.A201060.result}, of type \decl{Math0.Statements.A201060.statement}; the auxiliary declarations named in Section~\ref{sec:proof} are private to the proof file, except \decl{digits_rep}, \decl{isRepdigit_iff_mem_repList}, \decl{infinite_isTerm} and \decl{a_data}, which come from the statement file (together with \decl{mem_repList}). The residue tables and the end of the proof file read, verbatim (docstrings, blank lines and all other declarations omitted):
\begin{lstlisting}
private def T3 : List ℕ :=
  [1, 3, 9, 27, 81, 43, 29, 87, 61, 83, 49, 47, 41, 23, 69, 7, 21, 63, 89, 67]
private def T5 : List ℕ := [1, 5, 25]
theorem result : statement := by
  intro n hn
  have hterm : IsTerm (a n) := Nat.nth_mem_of_infinite infinite_isTerm (n - 1)
  have hge : 1111 ≤ a n := by
    rw [← a_sixteen]
    exact Nat.nth_monotone infinite_isTerm (by omega)
  exact isRepunit_of_isTerm hterm hge
end Math0.Frontier.A201060
#math0_settles Math0.Frontier.A201060.result Math0.Statements.A201060.statement
\end{lstlisting}
The proof file contains no \texttt{sorry}, \texttt{native\_decide}, \texttt{axiom} or \texttt{set\_option}; its finite checks (Lemma~\ref{lem:tables}) use \texttt{decide}, and the data lemmas of the statement file use \texttt{decide +kernel}. Both evaluate decision procedures in the Lean kernel and add no axioms.

\subsection*{Audit}
The command \texttt{\#math0\_settles}, defined in \path{lean/Math0/Audit.lean} for statements that math0 formalizes itself, fails unless (1) \decl{statement} is a \texttt{Prop}-valued definition without universe parameters; (2) \decl{statement} and \decl{result} use no axioms besides \decl{propext}, \decl{Classical.choice} and \decl{Quot.sound} (rejecting \decl{sorryAx}, the axioms introduced by \texttt{native\_decide}, and custom axioms); and (3) the type of \decl{result} is syntactically the constant \decl{statement} or its negation, up to metadata, with no unfolding. Here it reports
\begin{lstlisting}[basicstyle=\ttfamily\scriptsize]
math0 audit: 'Math0.Statements.A201060.statement' depends on axioms [Classical.choice, Quot.sound, propext]
math0 audit: 'Math0.Frontier.A201060.result' depends on axioms [Classical.choice, Quot.sound, propext]
math0 statement check: 'Math0.Frontier.A201060.result' proves 'Math0.Statements.A201060.statement'
\end{lstlisting}

\subsection*{Toolchain and reproduction}
We use Lean~4 \cite{Lean4} \texttt{v4.33.1} and Mathlib \cite{Mathlib} \texttt{v4.33.1}, pinned in \path{lean/}; the repository also pins formal-conjectures \cite{FC} at commit \texttt{c77fdd4}, which this result does not use. To reproduce, run \path{scripts/verify.sh} from the repository root: it fetches the Mathlib build cache, builds all math0 results with \texttt{lake build} (re-running every audit command), checks that the pinned \texttt{formal\_conjectures} and \texttt{mathlib} checkouts are unmodified, and reports the number of passed statement checks. With the dependencies built, the single file can be checked with \texttt{lake env lean Math0/Frontier/OEIS/A201060.lean} in the directory \path{lean}; this prints the three audit lines above.

{\small
\begin{thebibliography}{99}

\bibitem{BM99}
Y.~Bugeaud and M.~Mignotte, \emph{On integers with identical digits}, Mathematika \textbf{46} (1999), 411--417.

\bibitem{FC}
M.~Firsching, P.~Lezeau, S.~Mercuri, M.~Z. Horv\'ath, Y.~Dillies, C.~S\"onne, E.~Wieser, F.~Zhang, T.~Hubert, B.~Ag\"uera y Arcas, and P.~Kohli, \emph{Formal Conjectures: An Open and Evolving Benchmark for Verified Discovery in Mathematics}, arXiv:2605.13171, 2026. Repository: \url{https://github.com/google-deepmind/formal-conjectures}.

\bibitem{GP03}
A.~Gica and L.~Panaitopol, \emph{On Obl\'ath's problem}, J. Integer Seq.\ \textbf{6} (2003), Article~03.3.5.

\bibitem{KL05}
O.~Kihel and F.~Luca, \emph{Perfect powers with all equal digits but one}, J. Integer Seq.\ \textbf{8} (2005), Article~05.5.7, \url{https://cs.uwaterloo.ca/journals/JIS/VOL8/Kihel/kihel7.html}.

\bibitem{Mathlib}
The mathlib Community, \emph{The Lean mathematical library}, in: Proceedings of the 9th ACM SIGPLAN International Conference on Certified Programs and Proofs (CPP 2020), ACM, 2020, 367--381.

\bibitem{Lean4}
L.~de Moura and S.~Ullrich, \emph{The Lean 4 theorem prover and programming language}, in: Automated Deduction -- CADE 28, Lecture Notes in Comput.\ Sci.\ \textbf{12699}, Springer, 2021, 625--635.

\bibitem{Obl56}
R.~Obl\'ath, \emph{Une propri\'et\'e des puissances parfaites}, Mathesis \textbf{65} (1956), 356--364.

\bibitem{A002275}
The OEIS Foundation Inc., \url{https://oeis.org/A002275}, \emph{Repunits: $(10^n-1)/9$}, entry by N.~J.~A. Sloane.

\bibitem{A007954}
The OEIS Foundation Inc., \url{https://oeis.org/A007954}, \emph{Product of decimal digits of $n$}, entry by R.~Muller.

\bibitem{A010785}
The OEIS Foundation Inc., \url{https://oeis.org/A010785}, \emph{Repdigit numbers, or numbers whose digits are all equal}, entry by N.~J.~A. Sloane.

\bibitem{A201060}
The OEIS Foundation Inc., \url{https://oeis.org/A201060}, \emph{Repdigit numbers with repdigit digital product}, entry by J.~Krizek (Nov.~26, 2011); revision \#11, Mar.~30, 2012.

\end{thebibliography}
}
\end{document}
